\documentclass[a4paper]{article}

% Basic packages
% Español (México) con XeLaTeX: fontspec para fuentes Unicode, polyglossia
% para la división silábica y los nombres del documento en la variante mexicana.
\usepackage{fontspec}
\usepackage{polyglossia}
\setdefaultlanguage[variant=mexican]{spanish}
% Los nombres chinos van como «pinyin (original)»: xeCJK da a los caracteres
% Han su propia fuente; sin ella XeLaTeX los omite con un aviso.
% FandolSong no cubre algunos caracteres de nombres (U+73FA); WenQuanYi Zen Hei sí.
\usepackage[AutoFallBack=true]{xeCJK}
\setCJKmainfont{FandolSong-Regular.otf}
\setCJKfallbackfamilyfont{\CJKrmdefault}{WenQuanYi Zen Hei}
% polyglossia no traduce los nombres de listings.
\AtBeginDocument{\providecommand{\lstlistingname}{}\renewcommand{\lstlistingname}{Listado}%
  \providecommand{\lstlistlistingname}{}\renewcommand{\lstlistlistingname}{Índice de listados}}
\usepackage{amsmath, amssymb}
\usepackage{graphicx}
\usepackage[margin=2.5cm]{geometry}
\usepackage[most]{tcolorbox}
\usepackage{etoolbox}
\usepackage{listings}
\usepackage{hyperref}
\usepackage{booktabs}
\usepackage{subcaption}
\usepackage{float}

% Highlight boxes
\newtcolorbox{knowledgebox}[1]{
    enhanced, colback=blue!5!white, colframe=blue!75!black, colbacktitle=blue!75!black,
    coltitle=white, fonttitle=\bfseries, title={#1}, attach boxed title to top left={yshift=-2mm, xshift=2mm},
    boxrule=1pt, sharp corners
}

\newtcolorbox{importantbox}[1]{
    enhanced, colback=yellow!10!white, colframe=yellow!80!black, colbacktitle=yellow!80!black,
    coltitle=black, fonttitle=\bfseries, title={#1}, sharp corners
}

\newtcolorbox{warningbox}[1]{
    enhanced, colback=red!5!white, colframe=red!75!black, colbacktitle=red!75!black,
    coltitle=white, fonttitle=\bfseries, title={#1}, sharp corners
}

% Listings
\lstset{
    language=Python,
    basicstyle=\ttfamily\small,
    keywordstyle=\color{blue},
    stringstyle=\color{red!60!black},
    commentstyle=\color{green!60!black},
    breaklines=true,
    frame=single,
    numbers=left,
    numberstyle=\tiny\color{gray},
    captionpos=b,
    extendedchars=true
}

% Front-page metadata
\newcommand{\notetitle}{Zhang Bo (张钹), académico: Hacia la inteligencia artificial general}
\newcommand{\noteauthors}{Elaboradas a partir de materiales públicos del curso}
\newcommand{\notedate}{2025}
\newcommand{\videochannel}{AITIME Luntao (论道)}
\newcommand{\videopublishdate}{2025}
\newcommand{\videoduration}{aproximadamente 20 minutos}
\newcommand{\videourl}{}
\newcommand{\videocoverpath}{cover.jpg}
\newcommand{\repourl}{https://github.com/hqhq1025/ai-course-notes}


% Footer with repo info
\usepackage{fancyhdr}
\pagestyle{fancy}
\fancyhf{}
\fancyfoot[C]{\thepage}
\fancyfoot[R]{\footnotesize\color{gray}\href{\repourl}{AI Course Notes}}
\renewcommand{\headrulewidth}{0pt}
\renewcommand{\footrulewidth}{0pt}

\begin{document}

\begin{titlepage}
\centering
{\Large Notas de la charla\par}
\vspace{1.2cm}
{\huge\bfseries \notetitle\par}
\vspace{0.8cm}
{\large \noteauthors\par}
\vspace{0.3cm}
{\large \notedate\par}
\vspace{1.2cm}

\ifdefempty{\videocoverpath}{
{\small Al generar las notas, indique la ruta local de la portada del video.\par}
}{
\includegraphics[width=0.82\textwidth,height=0.45\textheight,keepaspectratio]{\videocoverpath}\par
}

\vfill
\begin{tcolorbox}[width=0.9\textwidth, colback=black!2!white, colframe=black!60, sharp corners]
\textbf{Autor o canal del video}: \videochannel\par
\textbf{Fecha de publicación}: \videopublishdate\par
\textbf{Duración del video}: \videoduration\par
\ifdefempty{\videourl}{}{
\textbf{Enlace del video}: \href{\videourl}{\nolinkurl{\videourl}}\par
}
\par
\textbf{Página del proyecto}: \href{\repourl}{\nolinkurl{\repourl}}\par
\end{tcolorbox}
\end{titlepage}

\tableofcontents
\newpage

%% --- Inicio del contenido --- %%

\section{Introducción: ¿qué domina en realidad un modelo de lenguaje grande?}

Zhang Bo (张钹), académico (decano honorario del Instituto de Inteligencia Artificial de la Universidad de Tsinghua, académico de la Academia de Ciencias de China), analizó de manera sistemática, desde la lingüística y la filosofía, los principios y las limitaciones de los modelos de lenguaje grandes, así como el verdadero reto de avanzar hacia la inteligencia artificial general (AGI).

Parte de una pregunta fundamental: ¿puede una máquina entender el lenguaje como lo hace un ser humano?

\begin{importantbox}{Logros y limitaciones de los modelos de lenguaje grandes}
Bajo indicaciones externas, las máquinas ya son capaces de generar, en dominios abiertos, un lenguaje diverso, coherente en el plano semántico y similar al humano. ¿Esto significa que dominan el lenguaje humano? La respuesta es: sí, pero \textbf{todavía no de manera completa}. En muchos aspectos, sigue habiendo una diferencia esencial con el lenguaje humano.
\end{importantbox}


\section{Semántica distribucional: los fundamentos teóricos de los modelos grandes}

\subsection{De lo disperso a lo denso}

El principio central de los modelos grandes de lenguaje es la \textbf{semántica distribucional}: usar el contexto que rodea a una palabra para definir el significado de esa palabra (la idea clásica de Firth).

\begin{knowledgebox}{Avance técnico clave}
Convertir el lenguaje de una \textbf{representación dispersa} en un espacio de alta dimensión a una \textbf{estructura geométrica de espacio vectorial denso} es un avance mayor. Con esto, el procesamiento del lenguaje se convierte por completo en un \textbf{problema de cálculo matemático}: el espacio de coocurrencia disperso original no era computable, pero al convertirse en un espacio vectorial denso puede procesarse con herramientas matemáticas.
\end{knowledgebox}

\subsection{Propiedades límite demostrables}

Se puede demostrar que, si la cantidad de datos es suficiente y el contexto es suficientemente largo, este espacio vectorial se aproxima al \textbf{espacio de relaciones semánticas}. Esto explica por qué la investigación actual persigue contextos más largos y más datos.

\begin{importantbox}{El significado de la aproximación semántica}
Una vez que el espacio vectorial se aproxima a las relaciones semánticas del lenguaje, en cierto sentido puede lograrse la ``comprensión'', e incluso puede lograrse la \textbf{reflexividad}: pensar sobre el propio pensamiento. Esto ya se ha observado en los modelos grandes de lenguaje.
\end{importantbox}

\subsection{Resumen de la sección}

Los modelos grandes de lenguaje se basan en el principio de la semántica distribucional y, mediante la conversión del espacio vectorial de disperso a denso, logran hacer computable el lenguaje. Pero esta definición de significado es en sí misma aproximada.

\section{Los cinco grandes defectos de la aproximación de los modelos}

\subsection{El dilema fundamental de la definición del significado}

El académico Zhang Bo (张钹) planteó una observación profunda: muchos fenómenos extraños que aparecen en los modelos de lenguaje grandes actuales \textbf{no pueden atribuirse simplemente a un problema de datos}.

\begin{warningbox}{Una atribución ampliamente malentendida}
«Mucha gente ve que en las máquinas aparecen muchos fenómenos extraños y lo atribuye a un problema de datos. Esto es completamente erróneo. Ahora mismo, muchas de estas cosas se deben a la \textbf{aproximación de los modelos}.» En el mundo existen siete u ocho escuelas filosóficas distintas sobre la definición del «significado», y ninguna de ellas es científicamente completa. La semántica distribucional no es más que una de esas aproximaciones.
\end{warningbox}

\subsection{Los cinco grandes defectos}

Debido a la incompletitud de la definición del significado, los modelos de lenguaje grandes presentan necesariamente los siguientes cinco defectos:

\begin{importantbox}{Los cinco defectos inherentes de los modelos de lenguaje grandes}
\begin{enumerate}
\item \textbf{Defecto de referencia}: el modelo no puede apuntar realmente a entidades del mundo real
\item \textbf{Defecto de veracidad y causalidad}: el modelo carece de capacidad de razonamiento causal y solo puede captar correlaciones estadísticas
\item \textbf{Defecto pragmático}: el modelo no puede comprender realmente las reglas de uso del lenguaje en contextos sociales específicos
\item \textbf{Defecto de polisemia y dinamismo}: la diversidad y el cambio dinámico del significado de las palabras rebasan la capacidad de una representación vectorial estática
\item \textbf{Defecto de contexto y comportamiento de bucle cerrado}: el modelo carece de capacidad de interacción de bucle cerrado con el entorno físico
\end{enumerate}
\end{importantbox}

El académico Zhang Bo (张钹) enfatizó: estos cinco defectos son \textbf{inherentes y necesarios} — no pueden eliminarse con más datos ni con modelos más grandes, porque la raíz del problema está en la incompletitud de la propia definición del significado.

\section{Definición operacional de la AGI: cinco capacidades clave}

\subsection{Por qué se necesita una nueva definición}

El académico Zhang Bo (张钹) señala que las definiciones actuales de la AGI (como la de Musk, «completar más del 70\% de las tareas humanas») \textbf{son completamente inejecutables e inverificables}, lo que provoca controversias inevitables.

\begin{warningbox}{El daño de las definiciones ambiguas}
Tomemos como ejemplo el reconocimiento de patrones: ¿basta con superar la tasa de reconocimiento humana para considerar que se ``alcanzó el nivel humano''? Unos dicen que sí, otros dicen que, en términos de robustez, aún falta mucho. Una definición ambigua hace que la discusión sobre la AGI caiga en una polémica interminable.
\end{warningbox}

\subsection{Cinco capacidades clave}

El académico Zhang Bo propone una definición de la AGI \textbf{ejecutable y verificable}, es decir, que debe alcanzar las siguientes cinco capacidades clave (presta especial atención al \textbf{adjetivo calificativo} de cada una):

\begin{importantbox}{Las cinco capacidades clave de la AGI}
\begin{enumerate}
\item Comprensión multimodal \textbf{consistente en espacio y tiempo}: la clave está en la ``consistencia espaciotemporal''; los ritmos temporales de las distintas modalidades no están sincronizados (el video avanza fotograma a fotograma, el texto puede resumir milenios en una oración), lo que hace la alineación extremadamente difícil
\item Aprendizaje y adaptación \textbf{en línea}: se distingue del aprendizaje fuera de línea; el reto central es la \textbf{controlabilidad}, es decir, si el proceso de aprendizaje por refuerzo puede converger hacia el objetivo
\item Razonamiento \textbf{verificable} y planeación \textbf{de largo plazo}: los resultados del razonamiento deben ser verificables; la planeación debe tener carácter de largo plazo
\item Reflexión y metacognición \textbf{verificables}: la reflexión actual es toda ``de sensación'', carente de una señal rastreable y verificable
\item Generalización \textbf{entre tareas}: no solo generalización entre dominios (algo que los modelos grandes ya hacen relativamente bien), sino generalización entre tareas fuera de distribución, con estructuras distintas y en escenarios de cola larga
\end{enumerate}
\end{importantbox}

\begin{knowledgebox}{Por qué importan los adjetivos}
El académico Zhang Bo insiste repetidamente en que ``los adjetivos son muy importantes''. Por ejemplo: muchos equipos trabajan en la ``comprensión multimodal'', pero al agregar ``consistente en espacio y tiempo'' la dificultad aumenta drásticamente; muchos modelos tienen ``razonamiento'', pero el razonamiento ``verificable'' es un requisito completamente distinto. Estos calificativos delimitan la verdadera frontera técnica.
\end{knowledgebox}

\section{De LLM a agente: rutas de implementación}

\subsection{Las seis grandes líneas técnicas actuales}

El académico Zhang Bo (张钹) considera que las seis cosas en las que la industria trabaja actualmente apuntan a los cinco objetivos mencionados arriba:

\begin{enumerate}
\item Fusión multimodal e interacción incorporada
\item Recuperación aumentada (RAG)
\item Alineación de conocimiento estructurado
\item Vinculación de herramientas y ejecución
\item Alineación y restricciones
\item Aprendizaje y adaptación en línea
\end{enumerate}

\subsection{Resumen de la sección}

Para pasar de un LLM a un agente capaz de ejecutar tareas complejas se necesitan avances en cinco dimensiones. Las seis líneas técnicas actuales corresponden cada una a algo, pero cada línea enfrenta el reto clave de pasar de «tenerlo» a «hacerlo bien».

\section{Los tres niveles de la agencia de la IA}

El académico Zhang Bo (张钹) propuso tres niveles de desarrollo de la inteligencia artificial como «agente»:

\begin{importantbox}{Los tres niveles de agencia}
\begin{enumerate}
\item \textbf{Agente de acción} (Agent of Action): la máquina como herramienta que ejecuta instrucciones humanas — ya \textbf{se alcanzó}; lo deseamos mucho, porque nos ayuda
\item \textbf{Agente de norma y responsabilidad} (Agent of Norm): la máquina es capaz de asumir responsabilidad — \textbf{aún no se alcanza}; la dificultad técnica es alta, pero es una meta que se puede perseguir
\item \textbf{Agente de experiencia y consciencia} (Agent of Consciousness): la máquina posee consciencia — este es el nivel que genera más preocupación
\end{enumerate}
\end{importantbox}

Para quienes trabajan en la práctica, el académico Zhang Bo recomienda prestar atención a los dos primeros niveles, sin considerar prematuramente el tercero.


\section{Alineación y gobernanza: el objeto de la gobernanza son las personas, no las máquinas}

\subsection{La paradoja de la alineación}

\begin{warningbox}{¿Con qué se alinea realmente la alineación?}
¿Las máquinas tienen que alinearse necesariamente con los seres humanos? Los propios seres humanos tienen defectos como la codicia y el engaño, defectos que las máquinas originalmente no tenían. La moral humana no es el estándar más alto. Entonces, ¿con qué debería alinearse realmente la «alineación»? Esto merece una discusión más profunda.
\end{warningbox}

\subsection{Gobernar a los investigadores y a los usuarios}

El académico Zhang Bo (张钹) señaló claramente que \textbf{el objeto de gobernanza más importante no son las máquinas, sino los seres humanos}, en particular los investigadores y los usuarios. Los empresarios de la era de la IA deben asumir una responsabilidad social mayor.

\section{La nueva misión de los empresarios en la era de la IA}

El académico Zhang Bo (张钹) no aprobaba antes que sus estudiantes emprendieran, pero la aparición de los modelos grandes cambió su opinión: \textbf{los mejores estudiantes deberían fundar empresas}, porque la inteligencia artificial redefinió el papel del empresario.

\begin{importantbox}{Seis nuevas responsabilidades}
El empresario de la era de la IA ya no se limita a ganar dinero, sino que debe:
\begin{enumerate}
\item \textbf{Redefinir la creación de valor}: no ofrecer simplemente productos y servicios, sino convertir el conocimiento y el razonamiento en herramientas reutilizables
\item Entregar la IA a la humanidad como una tecnología de uso general, igual que el agua y la electricidad
\item Asumir la responsabilidad social y las funciones de gobernanza
\item Lograr un crecimiento sostenible e incluyente
\item Participar en la construcción de la seguridad y la ética de la IA
\item Impulsar la educación y la formación de talento
\end{enumerate}
\end{importantbox}

\section{Síntesis y ampliación}

La ponencia del académico Zhang Bo (张钹) parte de fundamentos filosóficos y científicos, y ofrece un marco profundo para entender la ruta hacia la AGI:

\begin{enumerate}
\item \textbf{Fundamento teórico}: la semántica distribucional hace computable el lenguaje, pero la definición misma de significado es una aproximación
\item \textbf{Cinco limitaciones inherentes}: referencia, causalidad, pragmática, polisemia/dinamismo y conducta de circuito cerrado, que se originan en la aproximación del modelo y no en la falta de datos
\item \textbf{Definición operacional de la AGI}: cinco capacidades clave, y el adjetivo calificativo de cada una delimita la verdadera frontera tecnológica
\item \textbf{De los LLM a los Agentes}: seis direcciones técnicas que corresponden a cinco objetivos
\item \textbf{Tres niveles de la subjetividad}: sujeto de conducta $\rightarrow$ sujeto normativo $\rightarrow$ sujeto de conciencia
\item \textbf{El núcleo de la gobernanza}: gobernar a las personas (investigadores y usuarios), no a la máquina
\item \textbf{La nueva misión del empresario}: de ganar dinero a beneficiar a la humanidad; la IA redefine el papel del empresario
\end{enumerate}

\begin{knowledgebox}{Lecturas adicionales}
\begin{itemize}
\item Firth (1957): ``A Synopsis of Linguistic Theory''——fundamento filosófico de la semántica distribucional
\item Las siete grandes corrientes de la semántica (teoría referencial, teoría de las condiciones de verdad, pragmática, semántica cognitiva, etc.)
\item Definición de niveles de la AGI: el marco L1--L5 de OpenAI y la clasificación de AGI de Google DeepMind
\item Los avances de investigación más recientes en razonamiento verificable (Verifiable Reasoning)
\item Los planteamientos previos del académico Zhang Bo (张钹) sobre la ``tercera generación de inteligencia artificial''
\end{itemize}
\end{knowledgebox}

%% --- Fin del contenido --- %%

\end{document}
